Building the CR5 loop taught me that policy learning begins before model training. I anchored HDF5 records to RGB timestamps, aligned robot state, target action, and gripper feedback, then carried those fields through LeRobot conversion and WebSocket policy serving. I now want to study how synchronization and information flow shape sequential policy behavior across repeated experiments. CUHK's MSc in Information Engineering offers a structure for developing that inquiry from learning methods through successive research stages.

Subject to course availability and approval of the project sequence, I would combine IERG5350 Reinforcement Learning with IEMS5910 Advanced Research and Development Project I and IEMS5920 Advanced Research and Development Project II. The first project could compare synchronization choices in a sequential-learning pipeline; the second could examine how the resulting policy responds when sensing or serving conditions change. If a suitable placement were available, IEMS5930 Information Engineering Internship, lasting no less than twelve weeks, could expose the same methods to different operational constraints. My CR5 work provides a concrete starting case in which clock alignment and service handoffs change what a policy receives and executes. I want to turn that case into better information flows for embodied-AI deployment.
